\textbf{PART A:}

\begin{description}
\item[A.1] In cylindrical coordinates, the gravitational field is given by
  \[ \vec{g} = \vec{g}(r) = g(r)\,\hat{r} = g\,\hat{r}, \]
  \hfill[1 point]

  where, in view of the axial symmetry, $g$ can be computed from the Gauss
  law for the gravitational field:
  \[ \vec{g}\bullet\vec{A} = -\,4\pi G M_{in}, \]
  \hfill[1 point]

  where $M_{in}$ is the mass enclosed by the surface $A$, as shown in the
  figure below:
  % FIGURE NEEDED: hand-drawn Gaussian cylinder of length L and radius r
  % around the string (2021_T15_Solution.pdf page 1)

  Taking the surface to be a cylinder of length $L$ and radius $r$, from
  the above equation:
  \[ 2\pi r L g = -\,4\pi G M_{in} \]
  \[ g = -\frac{2GM_{in}}{rL} \]
  \hfill[1 point]

  There are two cases (regions):
  \begin{itemize}
  \item For $r > r_0$, one has $M_{in} = \mu L$ and so \hfill[0.5 points]
    \[ g = -\frac{2G\mu}{r} \]
    \hfill[1 point]
  \item For $r < r_0$, one has $M = \frac{r^2}{r_0^2}\mu L$ and so
    \hfill[0.5 points]
    \[ g = -\frac{2G\mu}{r_0}\left(\frac{r}{r_0}\right) \]
    \hfill[1 point]
  \end{itemize}

\item[A.2] Writing $\vec{g}(r_0)$, in terms of $G$, $\mu$ and $r_0$:
  \begin{align*}
  \vec{g}(r_0) &= -\frac{2G\mu}{r_0}\hat{r}\\
  g_0 \equiv \left|\vec{g}(r_0)\right| &= \frac{2G\mu}{r_0}
  \end{align*}
  \hfill[1 point]

\item[A.3] The following points should be included in the sketch:
  \begin{itemize}
  \item $(0, 0)$ \hfill[0.5 points]
  \item $(1, -1)$ \hfill[0.5 points]
  \end{itemize}
  Additionally,
  \begin{itemize}
  \item the sketch is linear between $(0,0)$ and $(1,-1)$
    \hfill[0.8 points]
  \item the sketch is concave for $r > r_0$,
    $\left(\sim\frac{1}{r}\right)$ \hfill[0.8 points]
  \item the sketch is drawn below x-axis. \hfill[0.4 points]
  \end{itemize}
  % FIGURE NEEDED: sketch of g(r): linear from (0,0) down to (r_0, -g_0),
  % then concave ~ -1/r rising back toward zero
  % (2021_T15_Solution.pdf page 2)

\item[A.4] The centripetal acceleration equals the gravitational
  acceleration, thus
  \[ \frac{v^2}{R} = \frac{2G\mu}{R} \]
  \hfill[1 point]
  \[ \frac{2\pi R}{\tau} = \sqrt{2G\mu} \]
  \hfill[1 point]
  \[ R = \frac{\sqrt{2G\mu}}{2\pi}\tau \]
  \[ A = \frac{\sqrt{2G\mu}}{2\pi} \]
  \hfill[1 point]
  \[ \alpha = 1 \]
  \hfill[1 point]

\item[A.5] The potential energy of the particle is given by
  \[ U(r) = -\int_b^r \vec{F}(r)\cdot\vec{dr} \]
  \[ = m\int_b^r g(r)\,dr \]
  \hfill[1 point]
  \[ = 2Gm\mu\int_b^r \frac{1}{r}\,dr \]
  \hfill[1 point]
  \[ = 2Gm\mu\,\ln\!\left(\frac{r}{b}\right) \]
  \hfill[1 point]

  Where $b$ is a constant that sets the 0 of $U$.

  Note 1: Usually, $U = 0$ is set at $b = \infty$. For a cosmic string,
  this is not possible. Instead, we can choose $U = 0$, for example, at
  the string's surface $b = r_0$ \emph{(this choice is not relevant!)}.
  Thus
  \[ U = 2Gm\mu\,\ln\!\left(\frac{r}{r_0}\right) \]

  Note 2: For completeness, notice that the potential inside the string,
  $r < r_0$ is
  \[ U = 2Gm\mu\left(\frac{r}{r_0}\right)^2 - U_0, \]
  where
  $U_0 = 2Gm\mu\left(\frac{r}{r_0}\right)^2
   - 2\,Gm\mu\,\ln\!\left(\frac{r}{r_0}\right)$
  is a constant (again, not relevant!) that can be chosen such that $U$
  is continuous at the surface of the string. However, for the solution
  of the problem, it is not necessary to show this.

\item[A.6] The total energy of the particle is conserved, thus:
  \[ \frac{1}{2}mv^2 + 2Gm\mu\,\ln\!\left(\frac{r}{b}\right)
     = 2Gm\mu\,\ln\!\left(\frac{R_{max}}{b}\right) \]
  \hfill[2 point]

  Solving for $R_{max}$:
  \[ R_{max} = Re^{\frac{v^2}{4G\mu}} \]
  \hfill[2 point]

  Notice that the answer does not depend on $b$.

\item[A.7] \textbf{NO.} \hfill[1 point]

  From the previous result,
  \[ R_{max} = Re^{\frac{v^2}{4G\mu}}, \]
  we see that it is \textbf{not possible} to escape the gravitational
  field since for any speed, there is always a maximum distance
  $R_{max} < \infty$
\end{description}

\textbf{PART B:}

\begin{description}
\item[B.1] The energy density is given by
  \[ \rho = aT^4 \]
  \hfill[1 point]

  Equivalently:
  \[ \rho = \frac{4\sigma}{c}T^4
     = \frac{\pi^2 k_B^4}{15\,\hbar^3 c^3}T^4 \]
  \hfill[1 point]

  Note 3: This result arrives from the integration of the spectral energy
  density given by the Planck law over all frequencies. However, it is
  not required to show this integral.

\item[B.2]
  \[ r_0 = \frac{\hbar^{n_1}c^{n_2}}{k_B T} \]
  \begin{itemize}
  \item $r_0$ has dimensions of length: $[l]$
  \item $\hbar$ has dimensions of energy $\times$ time: $[e\,t]$
  \item $c$ has dimensions of speed: $[l/t]$
  \item $k_B T$ has dimensions of energy: $[e]$
  \end{itemize}
  Then:
  \[ [l] = \frac{[e\,t]^{n_1}\left[\frac{l}{t}\right]^{n_2}}{[e]} \]

  Since the LHS and the RHS should have the same dimensions, we get:
  \begin{align*}
  [l]&:\ 1 = n_2\\
  [e]&:\ 0 = n_1 - 1\\
  [t]&:\ n_1 = n_2
  \end{align*}
  \hfill[2 points]

  The solution is
  \[ n_1 = 1 \]
  \hfill[1 point]
  \[ n_2 = 1 \]
  \hfill[1 point]

\item[B.3] The energy of a piece of the string of length $L$ is
  \[ E = \rho\pi r_0^2 L = Mc^2 \]
  \hfill[1 point]

  where $M = \mu L$ is the mass of the piece. Solving for $\mu$:
  \[ \mu = \frac{\rho\pi r_0^2}{c^2} \]
  \hfill[1 point]

\item[B.4] The weak field condition is
  \[ \frac{2G\mu}{c^2}\ll 1 \]
  \[ 2G\frac{\rho\pi r_0^2}{c^2}\ll 1 \]
  \hfill[1 point]
  \[ 2G\frac{\pi}{c^2}\left(aT^4\right)
     \left(\frac{\hbar c}{k_B T}\right)^2\ll 1 \]
  \hfill[1 point]
  \[ \frac{2Ga\hbar^2\pi}{c^2 k_B^2}\,T^2\ll 1 \]
  \hfill[2 point]

  Where we used: $\mu = \frac{\rho\pi r_0^2}{c^2}$, $\rho = aT^4$,
  $r_0 = \frac{\hbar c}{k_B T}$
  \begin{align*}
  \frac{2Ga\hbar^2\pi}{c^2 k_B^2}\,T^2 &\ll 1\\
  \frac{2G\hbar^2\pi}{c^2 k_B^2}
    \left(\frac{\pi^2 k_B^4}{15\,\hbar^3 c^3}\right)T^2 &\ll 1\\
  \frac{2\pi^3}{15}\left(\frac{Gk_B^2}{\hbar c^5}\right)T^2 &\ll 1\\
  \frac{2\pi^3}{15}\,\frac{T^2}{T_{Pl}^2} &\ll 1
  \end{align*}
  \hfill[2 point]

  where $T_{Pl} = 1.416784\times 10^{32}\,K$ (known as the Planck
  Temperature)

  Note 4: the numerical factor $\frac{2\pi^3}{15}\sim 4.13$,
  \emph{(a 5\% error tolerance is accepted)}

  Note 5: From \textbf{A.4} we get that
  \[ 2G\mu = v^2 \]
  where $v$ is the speed with which a particle would orbit a string. Thus
  the weak field is rewritten as
  \[ \frac{v^2}{c^2}\ll 1 \]
  which is equivalent to a non-relativistic condition.

\item[B.5] The weak field condition is
  \[ 4.13\left(\frac{T}{T_{Pl}}\right)^2\ll 1 \]
  equivalently
  \[ \sim 2\frac{T}{T_{Pl}}\ll 1 \]
  \hfill[1 point]
  \begin{description}
  \item[i.] $\frac{2T_{EW}}{T_{Pl}}\sim
    2\times\frac{10^{15}}{10^{32}}\approx 1.4\times 10^{-17}\ll 1$
    \hfill[1 point]
  \item[ii.] $\frac{2T_{GUT}}{T_{Pl}}\sim
    2\times\frac{10^{29}}{10^{32}}\approx 1.4\times 10^{-3}\ll 1$
    \hfill[1 point]
  \end{description}

  Note 6: Using the following expression
  $4.13\left(\frac{T}{T_{Pl}}\right)^2\ll 1$

  One gets
  \begin{description}
  \item[i.] $4.13\left(\frac{T_{EW}}{T_{Pl}}\right)^2
    \sim 2.1\times 10^{-34}\ll 1$
  \item[ii.] $4.13\left(\frac{T_{GUT}}{T_{Pl}}\right)^2
    \sim 2.1\times 10^{-6}\ll 1$
  \end{description}
  And are also acceptable answers \emph{(a 5\% error tolerance is
  accepted)}

\item[B.6]\
  \begin{description}
  \item[i.] Yes \hfill[0.5 point]
  \item[ii.] Yes \hfill[0.5 point]
  \end{description}
\end{description}

\textbf{PART C:}

\begin{description}
\item[C.1] From the figure, the asymptotic condition for the observer to
  see a second image is that the light ray travelling from O directly to
  S should bend and travel along SE. This is the maximum angle of
  bending.

  As all angles are small, we can safely use
  \[ \sin x \approx \tan x \approx x \]
  Now,
  \[ \frac{p}{SP} < \delta\phi \]
  \hfill[2 points]
  \begin{align*}
  &< \frac{4\pi G\mu}{c^2}\\
  &< 2\pi\left(\frac{2G\mu}{c^2}\right)\\
  &< 2\pi\left(\frac{2\pi^3}{15}\right)
     \left(\frac{T^2}{T_{Pl}^2}\right)
  \end{align*}
  \hfill[2 points]

  And $SP \approx \left(D_{OE} - D_{ES}\right)$
  \[ p < \frac{4\pi^4}{15T_{Pl}^2}\,T^2\left(D_{OE} - D_{ES}\right) \]
  \hfill[2 points]

  % FIGURE NEEDED: lensing geometry: points E, S, P on a line, observer O
  % above P; images O1, O2 along ES1, ES2; angles alpha, beta, delta-phi
  % (2021_T15_Solution.pdf page 7)

\item[C.2] In the figure, the blue lines represent the bending of two
  light rays corresponding to the images O1 and O2. Note that
  $D_{ES} \approx D_{ES1} \approx D_{ES2}$ and
  $D_{OS} \approx D_{OS1} \approx D_{OS2}$. Thus the angle
  $S1EO \approx S2EO \equiv \alpha$ and
  $S1OE \approx S2OE \equiv \beta$. $2\alpha$ is the angular separation
  we are looking for. \hfill[2 points]

  Further, notice
  \[ \alpha + \beta = \delta\phi \]
  \hfill[1 point]

  Using law of sines in the triangle \textbf{EOS1}
  \[ \frac{\alpha}{D_{S1O}} = \frac{\beta}{D_{ES1}} \]
  \hfill[1 point]

  we get
  \[ 2\alpha = 2\delta\phi
     \left(\frac{D_{OE} - D_{ES}}{D_{OE}}\right) \]
  \hfill[2 points]

\item[C.3] If $D_{OE} = 2D_{ES}$,
  \[ 2\alpha = \delta\phi
     = 2\pi\left(\frac{2\pi^3}{15}\right)
       \left(\frac{T_{GUT}^2}{T_{Pl}^2}\right)
     \approx 1.29\times 10^{-5} \]
  \hfill[2 points]

  and
  \[ \delta\phi = 1.22\frac{\lambda}{D} \]

  Thus, using
  $\lambda\in[3\times 10^{-7}m,\ 8\times 10^{-7}m]$
  \[ D = 1.22\frac{\lambda}{\delta\phi}
     \in\left[3.75\times 10^{-2}m,\ 7.51\times 10^{-2}m\right] \]
  \hfill[2 points]

  Any answer within this interval is valid.

  A remarkable result of this model is light deflection by a cosmic
  string, which leads to the possibility of detection through
  gravitational lensing. For instance, cosmic string trings moving across
  the line of sight with respect to an Earth-based observer will cause
  line-like discontinuities as shown in the figure below (taken from
  [1]).
  % sic: "cosmic string trings" as in source
  % FIGURE NEEDED: lensed spiral-galaxy images from Sazhin et al. (2007),
  % Fig. 6 reproduction (2021_T15_Solution.pdf page 9)

  [1] M. V. Sazhin, O. S. Khovanskaya, M. Capaccioli, G. Longo,
  M. Paolillo, G. Covone, N. A. Grogin, E. J. Schreier, Gravitational
  lensing by cosmic strings: what we learn from the CSL-1 case,
  \emph{Monthly Notices of the Royal Astronomical Society}, Volume 376,
  Issue 4, March 2007, Pages 1731--1739.
\end{description}
